% Transcription — fonds Alexandre Grothendieck, shelfmark 26, batch 1 (pages 1-10).
% Established from the facsimile archives/batches/26/batch-01.pdf (a mirror of
% https://grothendieck.umontpellier.fr/26.pdf), per the skill
% transcribe-grothendieck. The folder is ten pages and this is its only batch.
%
% Pages 1 and 9 are brown wrappers carrying, in his hand in the top right corner,
% "V / Plan / Notations" and "Problèmes ouverts" respectively, and nothing else:
% they take no \page{} and their words open the two sections of the file.
%
% The remaining eight pages are two distinct things.
% - Pages 2, 4, 6, 7, 8 and 10 are the notes the folder is named for: a
%   numbered plan of a chapter in four parts A) to D) (§1 to §22, pages 2 and
%   4, with the numbers of §8 to §11 and §18 to §21 overwritten as he reorders),
%   a two-line "Recommandation importante" and a question on equivalence
%   relations (page 6), a cancelled list of theorems with a diagram of twelve
%   letters A to L (page 7), a "Notations" sheet fixing the representable
%   functors Quot, Hilb, Div, Hom, Isom, Aut, Imm, the Weil restriction, Prépic
%   and Pic and their Hom/V formalism (page 8), and a list a) to j) of
%   "Problèmes-types à résoudre" with marginal cross-references to the plan's
%   sections and to a "Chap VI" (page 10). Page 7 is cancelled as a whole by
%   four long diagonal strokes, and its first five lines once more by short
%   ones; it is transcribed and the cancellation recorded once.
% - Pages 3 and 5 are the versos of the plan sheets 2 and 4: a carbon
%   typescript of his own, paginated 32 and 31 by the machine, on fibre
%   functors of the étale topos (its 7.9.1 to 7.9.3), with the Greek letters
%   and arrows inked by hand and the whole struck by one diagonal. It is his
%   mathematics, legible end to end, so it is transcribed; in the archivists'
%   order its page 32 comes before its page 31, and a note says so.
%
% Three readings are flagged rather than settled: the letter under the
% restriction symbol on pages 2, 4 and 8, an X overwritten (read Y after the
% Notations sheet); the name of the functor after "Div" on page 8, scribbled
% out; and the last word of the margin beside f) on page 10.
%
% Pass: Fable 5.1 (claude-fable-5-1), 2026-09-03 — first pass, unchecked against the pages by a human.
\documentclass[11pt,a4paper]{article}
% Le préambule commun à toutes les transcriptions du fonds Grothendieck.
%
% Il tient en une idée : **l'incertitude se compose, elle ne se résout pas.**
% Une transcription qui lisse un mot douteux a détruit la seule chose pour
% laquelle on la faisait — un lecteur doit pouvoir savoir, à chaque endroit,
% ce qui était sur le feuillet et ce qui a été ajouté. Les six macros qui
% suivent sont donc l'appareil critique, et non de la décoration.
%
% Les mêmes commandes sont reconnues par `scripts/render.mjs`, qui produit la
% vue de lecture du site. Une macro ajoutée ici doit l'être aussi là-bas,
% faute de quoi le rendu échouera bruyamment — ce qui est voulu : un rendu
% silencieusement faux est pire qu'un rendu absent.

\ProvidesPackage{grothendieck}[2026/08/08 Transcriptions du fonds Grothendieck]

\usepackage{amsmath,amssymb,amsthm}
\usepackage{mathrsfs}
\usepackage{stmaryrd}
% Les diagrammes commutatifs sont partout dans ce fonds : c'est la langue
% même de la géométrie algébrique de Grothendieck, et une transcription qui
% les rendrait en prose ne transcrirait rien.
\usepackage{tikz-cd}
% Français par défaut — c'est la langue des feuillets — mais l'anglais est
% chargé aussi : la traduction et le résumé sont des documents anglais, et la
% typographie française (espaces avant les deux-points) y serait une faute.
% Ces documents basculent par \selectlanguage{english} en tête de corps.
\usepackage[english,french]{babel}
\usepackage{fontspec}
\usepackage{geometry}
\geometry{a4paper,margin=25mm,marginparwidth=28mm}
\usepackage{marginnote}
\usepackage{xcolor}
% ulem, not soul. `soul`'s \st and \ul are built on a character-by-character
% scan that XeLaTeX's font machinery breaks: the document fails with an
% undefined control sequence rather than degrading. ulem does the same two jobs
% and is Unicode-engine safe. `normalem` stops it hijacking \emph, which the
% transcriptions use for Grothendieck's own emphasis.
\usepackage[normalem]{ulem}
\usepackage{hyperref}
\hypersetup{colorlinks=true,linkcolor=black,urlcolor=black,pdfborder={0 0 0}}

% --- Métadonnées du lot -----------------------------------------------------
% Reprises telles quelles par le script de rendu, qui les affiche en tête de
% la vue de lecture. Elles fixent ce que le fichier prétend couvrir : sans
% elles, une transcription flotte sans point d'attache dans un fonds de seize
% mille feuillets.

\newcommand{\folder}[1]{\gdef\grothFolder{#1}}
\newcommand{\batch}[1]{\gdef\grothBatch{#1}}
\newcommand{\leaves}[2]{\gdef\grothFirst{#1}\gdef\grothLast{#2}}
\newcommand{\foldertitle}[1]{\gdef\grothFoldertitle{#1}}
\gdef\grothFolder{?}\gdef\grothBatch{?}\gdef\grothFirst{?}\gdef\grothLast{?}\gdef\grothFoldertitle{}

% --- Le résumé --------------------------------------------------------------
%
% Il ouvre la lecture modernisée, sous le titre « Résumé », et il est composé
% en petit. Ce n'est pas un ornement : c'est le seul passage du document qui
% ne soit pas des mathématiques, et un lecteur doit pouvoir le reconnaître
% avant de l'avoir lu — puis le sauter s'il connaît déjà le sujet, ou s'y
% arrêter s'il ne le connaît pas. Le filet qui le suit marque la reprise du
% corps.

\newenvironment{resume}
  {\small\setlength{\parskip}{0.45em}}
  {\par\medskip\hrule\medskip}

% --- Mots-clés ---------------------------------------------------------------
%
% En anglais, en clôture du résumé de la lecture modernisée : le vocabulaire
% moderne sous lequel chercher ce que les feuillets construisent. Le manifeste
% du site (`npm run manifest`) les extrait pour étiqueter la cote entière —
% cette ligne est la source unique des tags, il n'y a pas de fichier à côté.
\newcommand{\keywords}[1]{\par\smallskip\noindent
  {\footnotesize\textsc{Keywords} --- \textit{#1}}\par}

% --- L'appareil critique ----------------------------------------------------

% Le numéro de feuillet, en marge. C'est l'ancre qui permet de régler un
% désaccord sur une formule en regardant *un* feuillet plutôt qu'une chemise
% de six cents. Le site s'en sert aussi pour tourner les pages du fac-similé
% quand on fait défiler la transcription.
\newcommand{\leaf}[1]{%
  \marginnote{\footnotesize\color{gray!70}\textsf{#1}}%
  \label{leaf:#1}\ignorespaces}

% Illisible : on ne devine pas. Un mot inventé qui se lit comme les autres est
% le pire résultat possible.
\newcommand{\ill}{\textcolor{red!60!black}{[\,\ldots\,]}}

% Lecture douteuse : le mot est donné, et signalé comme incertain.
\newcommand{\uncertain}[1]{\uline{#1}}

% Ajout de l'éditeur — une lettre manquante, un mot sous-entendu.
\newcommand{\add}[1]{\textcolor{blue!50!black}{[#1]}}

% Biffé par Grothendieck lui-même. Ce qu'il a rayé fait partie du document :
% une rature dit souvent où la pensée a changé de direction.
\newcommand{\struck}[1]{\sout{#1}}

% Note du transcripteur, sur l'état matériel du feuillet.
\newcommand{\note}[1]{\par\smallskip{\small\color{gray!60!black}\textit{[#1]}}\par\smallskip}

% Note marginale de Grothendieck, distincte du corps du texte.
\newcommand{\marginal}[1]{\par\smallskip\noindent\hspace{1em}%
  {\small\color{gray!60!black}#1}\par\smallskip}

% --- Titre ------------------------------------------------------------------

\AtBeginDocument{%
  \begin{center}
    {\large\bfseries Fonds Alexandre Grothendieck --- cote n\textsuperscript{o}~\grothFolder}\\[2pt]
    {\itshape\grothFoldertitle}\\[4pt]
    {\small feuillets \grothFirst--\grothLast\ (lot \grothBatch)}\\[2pt]
    {\footnotesize\color{gray!60!black}Transcription établie d'après le fac-similé numérisé
     par l'Université de Montpellier.\\
     Les lectures incertaines sont soulignées, les illisibles notées [\,\ldots\,],
     les ajouts entre crochets.}
  \end{center}
  \vspace{1em}}

\endinput


\folder{26}
\batch{1}
\pages{1}{10}
\dating{[après 1967]}
\watermark{Édition de démonstration}
\foldertitle{EGA VI. Plans, notations, problèmes ouverts : notes manuscrites (s.d.)}

\begin{document}

\section*{V — Plan. Notations}

\note{titre de sa main, à l'encre, dans l'angle supérieur droit du premier
feuillet — une chemise brune qui ne porte rien d'autre : « V » souligné, puis
« Plan », « Notations ». Le feuillet ne reçoit pas de \texttt{page} ; ses mots
ouvrent la section}

\page{2}

\subsection*{A) Descente et passage au quotient}

§ 1 \quad Descente fidèlement plate

§ 2 \quad \struck{Passage aux quoti\ill{}} Topologies sur $(\mathrm{Sch})_{/S}$.
Propriétés générales des \struck{Généralités sur les foncteurs sur} foncteurs
sur $(\mathrm{Sch})_{/S}$. [Intervertir 1 et 2 ??] \struck{$(= \mathrm{Sch})_{/S}$}

\note{la ligne du § 2 est réécrite par-dessus une première rédaction biffée ;
« Topologies sur $(\mathrm{Sch})_{/S}$ » est interligné au-dessus, « foncteurs
sur $(\mathrm{Sch})_{/S}$ » au-dessous, d'une encre plus bleue}

§ 3 \quad Passage au quotient plat.

§ 4 \quad Descente non plate \quad [Applications : Ram. Samuel généralisé\ldots
\quad Appl. à platitude \uncertain{d'une relation schématique} et lemme de
Hironaka.

\note{la seconde ligne du crochet, « Appl. à platitude \ldots{} Hironaka », est
ajoutée en dessous d'une encre plus bleue ; le crochet n'est pas refermé}

§ 5 \quad Passage au quotient non plat. [Passage aux \uncertain{quots en
gén.}]

§ 6 \quad Descente de schémas relatifs étales, \ldots{} \quad [\uncertain{Évt}
descente des \struck{\ill{}} faisceaux étales ?? \ldots]

\subsection*{B) Méthodes non projectives de construction}

\note{titre reconstitué de ses corrections : « Propriétés générales des
foncteurs sur $(\mathrm{Sch})_{/S}$ » est biffé en entier ; « Méthodes non
projectives » est écrit en dessous, et « de construction », entouré, est
rattaché par un trait à la suite de « Méthodes ». Vient ensuite une ligne
encadrée puis biffée, numérotée d'un § lui-même biffé : « Propriétés générales
des foncteurs et propriétés de permanence »}

\struck{§ \ill{} \quad Propriétés générales des foncteurs et propriétés de
permanence}

§ 7 \quad Proreprésentabilité (et variantes).

§ 9 \quad Invariants différentiels \struck{de} tangents et normaux

\note{le numéro est un 8 surchargé en 9, avec une flèche au-dessus ; « tangents
et normaux » est interligné au-dessus du mot biffé}

§ 10 \quad (Après § 12 ??) \quad Propriétés générales de foncteurs, propriétés
de permanence
$\left\{\;\text{Foncteurs excellents sur } S \text{ ; morphismes excellents ??}\right.$

\note{« § 10 » est entouré, « Après § 12 ?? » encadré au-dessus de la ligne ;
l'accolade et son contenu sont dans la marge droite}

§ 11 \quad Représentabilités élémentaires du type $\prod_{Y/S} X/Y$. \quad
(Comme cas \ill{} ? Cas Greenberg ?) \quad cf SGAD VIII 6 et XI 6.8.

\note{le numéro est un 10 biffé, remplacé par 11 ; le tout entouré et relié par
un trait à « § 10 » et à « § 8 » — les trois numéros sont en cours de
redistribution. La lettre sous le $\prod$ est un X surchargé, lu Y d'après la
feuille de notations (p. 8) où la même surcharge se lit ; de même p. 4}

§ 8 \quad Construction de Modules cohérents

\note{le numéro, encadré, est surchargé et se lit 8 ; une première valeur
n'est plus lisible}

§ 12 \quad \struck{Représentabilité} Foncteurs représ. par morphismes loc.
pr. f.

\note{« Foncteurs » est interligné au-dessus du mot biffé, « représ. » au-dessus
de « repr. »}

§ 13 \quad Applications : des ths de représentabilité relatifs. (y compris
\uncertain{pour les} foncteurs étales\ldots)

\page{3}

\subsection*{Dactylographie, p. 32}

\note{pages 3 et 5 : dactylographie de sa main, paginée 32 et 31 par la
machine, lettres grecques et flèches encrées à la main, le tout biffé d'un
long trait diagonal. Dans l'ordre des archivistes sa page 32 (ici) précède sa
page 31 (p. 5), qui la précède dans le texte ; le paragraphe 7.9.3 ouvert au
bas de la p. 5 se poursuit ici}

Considérant la restriction de $\varphi$ à la sous-catégorie pleine de
$\widetilde{X}_{\text{ét}}$ formée des ouverts de ce topos, ou ce qui revient
au même (6.1.) des ouverts $U$ de $X$, on trouve une application
\[
  \varphi_0 : \mathcal{U}(X) \longrightarrow \{0, 1\}
\]
commutant aux sup quelconques et aux inf finis. (On note que pour $U \in
\mathcal{U}(X)$, $\varphi(U)$ est vide ou réduit un point, et on prend
$\varphi_0(U) = 0$ ou $1$ suivant qu'on est dans l'un ou l'autre cas). On voit
facilement \struck{\ill{}}, grâce au fait que tout fermé irréductible de $X$
a un point génerique et un seul, que $\varphi_0$ est defini à l'aide d'un
\struck{\ill{}} unique $x \in X$, par la condition
\[
  \varphi_0(U) = 1 \quad \text{sss} \quad x \in U ,
\]
i.e.
\[
  \varphi(U) \neq \emptyset \iff x \in U .
\]
Soit $F \in \mathrm{Ob}\,\widetilde{X}_{\text{ét}}$, alors l'image de $F$ dans le
faisceau final $X$ est un ouvert ; et comme
\[
  \varphi(F) \longrightarrow \varphi(U)
\]
est surjectif, on voit que $\varphi(F) \neq \emptyset$ sss $\varphi(U) \neq
\emptyset$ i.e. sss $x \in U$. Soit alors $C_\varphi$ la catégorie des couples
$(X', \xi)$ où $X' \in \mathrm{Ob}\,X_{\text{ét}}$, et $\xi \in \varphi(X')$. On a
signalé dans 7.9.1. que $C_\varphi^{\circ}$ est pseudo-filtrante connexe. De
plus grâce à 7.9.2. et à ce qui précéde appliqué à $X'$ au lieu de $X$, si
$(X', \xi) \in \mathrm{Ob}\,X_{\text{ét}}$, il lui est associé un unique point
$x' \in X'$, tel que pour un morphisme étale $X'' \longrightarrow X'$, $\xi$
est dans l'image de $\varphi(X'') \longrightarrow \varphi(X')$ sss $x'$ est
dans celle de $X''$. On conclut de ceci que le préschéma $\varprojlim$ des $X'$
suivant la catégorie $C$ des $(X', \xi)$ existe et est un localisé strict $Z$
de $X$, correspondant donc à un point géométrique $\xi$ et un foncteur fibre
\uncertain{$\mathcal{E}_{\xi}$}.

\note{« réduit un point », « génerique », « defini », « précéde » : sic. Les
s finals de « les préschémas » sont recouverts à la machine, d'où « le
préschéma ». Le dernier symbole est un glyphe manuscrit épais, indexé $\xi$,
qui n'est pas le $\xi$ des occurrences précédentes}

\page{4}

\subsection*{C) Méthodes projectives de construction}

\note{« de construction » est interligné au-dessus de « projectives », qui est
encadré, et rattaché à « Méthodes » par un trait}

§ 14 \quad Familles \uncertain{limitées}.

§ 15 \quad Foncteurs et Schémas de Hilbert. \quad (+ appl. aux quotients).

§ 16 \quad Foncteurs et Schémas de Picard. \quad (Ths d'existence projectifs
exclusivement ; également le cas de \uncertain{groupes structuraux affines}
commutatifs plus \uncertain{généraux} que $\mathbb{G}_m$\ldots)

\note{« Foncteurs et » est interligné au-dessus de « Schémas » aux § 15 et
§ 16, avec une parenthèse ouvrante ajoutée devant « Schémas ». La parenthèse
du § 16 court sur quatre lignes dans la marge droite, contre une accolade qui
regroupe les § 16 à § 20}

§ 17 \quad Variantes de Picard.

§ 18 \quad Applications aux courbes

§ 18 \quad Formalisme des normes

§ 20 \quad Invariants \uncertain{multiprojectifs} de Mumford.

\note{la colonne des numéros § 16 à § 20 est surchargée : le premier « 18 »
est écrit sur un autre chiffre, le second « 18 » et le « 20 » aussi, et la
numérotation ne se suit plus. Elle est transcrite telle qu'elle se lit}

\subsection*{D) Méthodes mixtes de construction}

\note{« de construction » interligné au-dessus de « mixtes », entre
parenthèses}

§ 21 \quad Cas de représentabilité de $\underline{\mathrm{Pic}}$,
$\underline{\mathrm{Pic}}^{\circ}$, $\underline{\mathrm{Pic}}^{\tau}_X$\ldots
\quad (y compris le cas de Mumford utilisant § 20).

\note{le numéro se lit « 20 » surchargé en « 21 » ; le « § 20 » de la
parenthèse est lui aussi surchargé}

§ 22 \quad Cas de représentabilité des foncteurs $\prod_{Y/S} X/Y$ etc.

\note{sous le $\prod$, un X surchargé, lu Y (cf. p. 2, § 11)}

(Ths du type Matsumura).

NB. \quad On suppose connus, tous les gén. fondamentaux (topologies, descente,
faisceaux, quotient) ; i.e. \uncertain{rappelés} au besoin dans
\uncertain{V}.

\note{le NB est isolé entre deux traits horizontaux ; le « V » est souligné et
précédé d'un petit rond, comme sur la chemise}

Manquent :

\struck{Il manquerait}

\begin{itemize}
  \item dans Chap. sur schémas en groupes : \{ les ths relatifs aux VA ; les ths
  genre finitude sur Picard, et l'étude qualitative (générale des schémas de
  Picard, (y compris critère de lissité de Mumford\ldots)
  \item dans Chapitre sur Picard : \{ (1.) \struck{Relation entre} Étude des
  Picard d'un schéma projectif en termes d'un pinceau de sections hyperplanes.
  Les th. sur Picard genre « \uncertain{certaines équivalences} ».
  \item dans Chap. sur $\pi_1$ : \{ Les formulations de descente : le van
  Kampen etc, en termes du groupe fondamental
\end{itemize}

\note{les trois rubriques sont écrites dans la marge gauche, chacune contre
une accolade qui enserre ses lignes ; « dans Chap. sur schémas en groupes »
est relié par un trait à l'accolade des deux premières lignes, sous « Il
manquerait » biffé. « qualitative » est interligné au-dessus de « (générale »}

\page{5}

\subsection*{Dactylographie, p. 31}

\note{la page commence au milieu d'une phrase ; ce qui précède n'est pas dans
le dossier}

(les morphismes se définissant de la façon évidente), et on note qu'on a un
homomorphisme évident
\[
  (\ast) \qquad \varinjlim_{(Y, \xi) \in C_\varphi^{\circ}} \check{Y}
  \longrightarrow \varphi
\]
en remarquant que
\[
  \mathrm{Hom}\bigl(\varinjlim_{C_\varphi} \check{Y}, \varphi\bigr) \simeq
  \varprojlim_{C_\varphi} \mathrm{Hom}(\check{Y}, \varphi) \simeq
  \varprojlim_{C_\varphi} \varphi(\widetilde{Y}) .
\]
Or on a un élément canonique dans $\varprojlim_{C_\varphi} \varphi(\widetilde{Y})$,
en associant à tout $(Y, \xi)$ l'élément $\xi$ de $\varphi(Y)$. Utilisant le
fait que les $\varprojlim$ finies existent dans $C$, et que le foncteur $Y
\rightsquigarrow \varphi(\widetilde{Y})$ y commute, on voit que dans $C_\varphi$
les $\varprojlim$ finies existent ; à fortion $C_\varphi^{\circ}$ est
pseudo-filtrant. Utilisant que tout faisceau $F$ est limite inductive des
faisceaux de la forme $\widetilde{Y}$, et utilisant le fait que $\varphi$
commute aux limites inductives, on conclut aisement que l'homomorphisme de
foncteurs $(\ast)$ est bijectif, et donne donc la représentation annoncée.

\note{« à fortion », « aisement » : sic}

7.9.2. \quad Remarquons également que si $\underline{E}$ est un topos, $Y$ un
objet de $\underline{E}$, d'où un topos $\underline{E}_{/Y}$, alors la donnée
d'un foncteur fibre $\varphi$ pour $\underline{E}_{/Y}$ équivaut à la donnée
d'un couple $(\varphi, \xi)$, où $\varphi$ est un foncteur fibre pour
$\underline{E}$, et $\xi \in \varphi(Y)$. Au foncteur fibre $\psi :
\underline{E}_{/Y} \longrightarrow (\mathrm{Ens})$ on associe le couple
$(\varphi, \xi)$, où $\varphi$ est le composé $\varphi(Z) = \psi(Z \times Y)$,
$\xi \in \varphi(Y) = \psi(Y \times Y)$ est l'image de l'unique élément de
$\psi(Y)$ par le morphisme diagonal de $Y \times Y$.

7.9.3. \quad Revenons maintenant au cas du site $C = \widetilde{X}_{\text{ét}}$,
et soit $\varphi : \widetilde{X}_{\text{ét}} \longrightarrow (\mathrm{Ens})$ un
foncteur fibre sur le topos étale de $X$.

\note{la suite est la p. 3}

\page{6}

\subsection*{Recommandation importante}

Partout où c'est possible, remplacer « fid. plt qu.-cpt » par « couvrant ».

\note{un trait horizontal sépare la recommandation de ce qui suit}

(pré)-Relations d'équivalence en pièces détachées !?
\[
  (R_0^{i}) \leftleftarrows (R_1^{ij})
\]

\note{la formule est entourée ; devant elle un signe biffé, illisible. Des deux
flèches entre $(R_1^{ij})$ et $(R_0^{i})$, seule l'inférieure porte une pointe
lisible}

\page{7}

\note{page biffée en entier par quatre longs traits diagonaux ; les cinq
premières lignes le sont une seconde fois par des traits courts. Dans l'angle
supérieur gauche, un gribouillis. Tout est transcrit}

Th. de finitude pour les schémas de Hilbert

Th. de Matsusaka pour les automorphismes d'une variété polarisée

Th. de la \struck{th} « courbe générique » \quad [pour Picard]

Schémas de Hilbert d'une courbe \struck{\ill{}} sur $S$.

Picard des courbes sur $S$ \quad [« diviseur $\Theta$ » etc ?]

Picard d'un schéma propre non projectif sur $k$

Picard d'un $S$-schéma projectif à fibres réductibles \quad [p.ex. cas d'une
courbe sur $S$]

\begin{tikzcd}[column sep=small, row sep=large]
  A \arrow[r] & B \arrow[d] \arrow[drr] & C \arrow[dl] \arrow[dr] & D \arrow[dr] \arrow[rr, bend left=25] & E \arrow[d] & J & K \\
  & H \arrow[dr] & & G \arrow[dl] & F \arrow[l] & & \\
  & & I & & & &
\end{tikzcd}

\note{les douze lettres sont posées sur trois rangs, sans autre légende ; la
flèche $A \to B$ est ondulée, les flèches $B \to G$, $C \to G$ et $D \to F$
sont tracées d'un trait épais, les autres d'un trait fin. Les traits diagonaux
qui biffent la page traversent aussi le diagramme}

$L$ \qquad Appendice : Pbs ouverts.

\note{sous un trait horizontal qui ferme le diagramme}

\page{8}

\subsection*{Notations}

\note{le titre est dans l'angle supérieur droit, souligné, séparé de la première
ligne par un trait vertical. Sur toute la hauteur de la marge droite, contre
un trait vertical, une colonne de lettres en regard des lignes : F, X, X, X,
X, X, X, Y, puis X, X, G. Chaque nom de foncteur est écrit trois fois sur sa
ligne : doublement souligné, simplement souligné, non souligné ; deux
flèches de sa main, « préfaisceau » et « faisceau », désignent les
occurrences simplement soulignées des lignes Prépic et Prépic$'$}

\marginal{F, X, X, X, X, X, X, Y, X, X, G}

\begin{align*}
  \underline{\underline{\mathrm{Quot}}}_{F/X/S}(S')
    &= \Gamma\bigl(S', \underline{\mathrm{Quot}}(F'/X'/S')\bigr)
     = \mathrm{Quot}(F'/X'/S') \\
  \underline{\underline{\mathrm{Hilb}}}_{X/S}(S')
    &= \Gamma\bigl(S', \underline{\mathrm{Hilb}}(X'/S')\bigr)
     = \mathrm{Hilb}(X'/S')
     \qquad = \underline{\underline{\mathrm{Quot}}}_{\mathcal{O}_X/X/S}(S') \\
  \underline{\underline{\mathrm{Div}}}_{X/S}(S')
    &= \Gamma\bigl(S', \underline{\mathrm{Div}}(X'/S')\bigr)
     = \mathrm{Div}(X'/S')
\end{align*}

\note{aux trois occurrences de « Div » le nom continue par des lettres
gribouillées jusqu'à l'illisible, comme s'il en avait essayé un plus long ;
au-dessus du premier, une apostrophe}

\begin{align*}
  \underline{\underline{\mathrm{Hom}}}_S(X, Y)(S')
    &= \Gamma\bigl(S', \underline{\mathrm{Hom}}_{S'}(X', Y')\bigr)
     = \mathrm{Hom}_{S'}(X', Y') = \mathrm{Hom}_S(X \times_S S', Y) \\
  \underline{\underline{\mathrm{Isom}}}_S(X, Y)(S')
    &= \Gamma\bigl(S', \underline{\mathrm{Isom}}_{S'}(X', Y')\bigr)
     = \mathrm{Isom}_{S'}(X', Y') \\
  \underline{\underline{\mathrm{Aut}}}_S(X)(S')
    &= \Gamma\bigl(S', \underline{\mathrm{Aut}}_{S'}(X')\bigr)
     = \mathrm{Aut}_{S'}(X') \\
  \underline{\underline{\mathrm{Imm}}}_S(X, Y)(S')
    &= \Gamma\bigl(S', \underline{\mathrm{Imm}}_{S'}(X', Y')\bigr)
     = \mathrm{Imm}_{S'}(X', Y')
\end{align*}

\note{à la ligne Aut, après $(X)$, un $(S')$ recouvert de traits est réécrit à
côté}

\[
  \Bigl(\underline{\underline{\prod}}_{Y/S} X/Y\Bigr)(S')
  = \Gamma\bigl(S', \underline{\Gamma}_{Y'/S'}(X'/Y')\bigr)
  = \bigl[\Gamma_{Y'/S'}(X'/Y') = \bigr]\;\Gamma(X'/Y')
\]

\note{dans le $\Gamma(S', \ldots)$, un mot biffé jusqu'à l'illisible précède
$\underline{\Gamma}_{Y'/S'}$. Sous $\underline{\underline{\prod}}_{Y/S} X/Y$, une
accolade et le mot « ou », puis, à la marge gauche :}

\[
  \underline{\underline{\Gamma}}_{Y/S}(X/Y)
  = \underline{\underline{H}}^{0}_{Y/S}(X/Y)
\]

\note{les deux barres de fraction de cette ligne sont repassées d'un trait
épais, et la ligne se termine par un signe biffé ; le Y sous le $\prod$ de la
ligne précédente est un X surchargé}

\begin{align*}
  \underline{\underline{\mathrm{Prépic}}}_{X/S}(S')
    &= \Gamma\bigl(S', \underline{\mathrm{Prépic}}_{X'/S'}\bigr)
     = \mathrm{Prépic}(X'/S')
     \qquad = H^1(X', \mathcal{O}_{X'}^{*}) \\
  \underline{\underline{\mathrm{Prépic}}}'_{X/S}(S')
    &= \Gamma\bigl(S', \underline{\mathrm{Prépic}}'_{X'/S'}\bigr)
     = \mathrm{Prépic}'(X'/S')
     \qquad = \Gamma\bigl(S', \underline{H}^{1}_{X'/S'}(\underline{\mathcal{O}}^{*}_{X'})\bigr) \\
  \underline{\underline{\mathrm{Pic}}}_{X/S}(S')
    &= \Gamma\bigl(S', \underline{\mathrm{Pic}}_{X'/S'}\bigr)
     = \mathrm{Pic}(X'/S')
     \qquad = \widetilde{\underline{\underline{\mathrm{Prépic}}}}_{X/S}(S')
\end{align*}

\note{à la ligne Prépic, l'indice de $\underline{\mathrm{Prépic}}$ est un pâté
sous lequel se lit $X'/S'$ ; la ligne continuait par « $= H^1(\ldots) =
\Gamma(S', H^1_{\ldots})$ », recouvert de hachures jusqu'à l'illisible, et
« $= H^1(X', \mathcal{O}_{X'}^{*})$ » est écrit au-dessus. À la ligne
Prépic$'$, une lettre est raturée dans le premier « Prépic », un $\neq$ biffé
suit $\mathrm{Prépic}'(X'/S')$, et « faisceau » est écrit sous
$\underline{H}^{1}_{X'/S'}(\underline{\mathcal{O}}^{*}_{X'})$. À la ligne Pic,
le tilde couvre « Prépic » et l'argument commence par un pâté}

2. \quad
\begin{align*}
  \mathrm{Hom}\Bigl(\underline{\underline{\prod}}_{X/S} \underline{F}, \underline{M}\Bigr)
    &= H^0\bigl(\underline{F} \otimes_S \underline{M}\bigr), &
  \underline{\mathrm{Hom}}\Bigl(\underline{\prod}_{X/S} \underline{F}, \underline{M}\Bigr)
    &= \underline{H}^{0}_{X/S}\bigl(\underline{F} \otimes_S \underline{M}\bigr) \\
  \underline{V}\Bigl(\prod_{X/S} \underline{F}\Bigr)
    &= \prod_{X/S} V(\underline{F}) &
  \underline{V}\Bigl(\underline{\underline{\prod}}_{X/S} \underline{F}\Bigr)
    &= \underline{\underline{\mathrm{Hom}}}_{X/S}(F, \underline{\mathcal{O}}_X)
\end{align*}

\note{au-dessus des deux $\prod_{X/S}\underline{F}$ de la première ligne,
« contravariant en F » (le premier suivi d'un mot biffé), avec une accolade.
La formule de droite de la seconde ligne est entourée, précédée d'un « ? », et
suivie dans la marge de « Notation opportune ??? »}

\struck{$\underline{\underline{\mathrm{Hom}}}_{X/S}(F, G)(S') = \Gamma\bigl(S',
\underline{\mathrm{Hom}}_{X'/S'}(F', G')\bigr) = \mathrm{Hom}_{\mathcal{O}_{X'}}(F', G')$}

\note{ligne encadrée puis biffée de hachures, avec au-dessus un second départ,
« $\underline{\underline{\mathrm{Hom}}}_{X/S}(F, G)(S')$ », biffé de même}

\begin{align*}
  \mathrm{Hom}\bigl(\underline{\underline{\mathrm{Hom}}}_{X/S}(F, G), \underline{M}\bigr)
    &= \mathrm{Hom}_{\mathcal{O}_X}(F, G \otimes_{\mathcal{O}_S} M) \\
  \mathrm{Hom}_{\mathcal{O}_{S'}}\bigl(\underline{\underline{\mathrm{Hom}}}_{X/S}(F, G) \otimes_{\mathcal{O}_S} S', \mathcal{O}_{S'}\bigr)
    &= \mathrm{Hom}_{\mathcal{O}_{X'}}(F', G') \\
  \underline{\mathrm{Hom}}_{\mathcal{O}_S}\bigl(\underline{\underline{\mathrm{Hom}}}_{X/S}(F, G), \underline{M}\bigr)
    &= \underline{\mathrm{Hom}}_{X/S}(F, G \otimes M)
\end{align*}

\note{les deux premières lignes sont à gauche, prises dans une accolade, la
troisième à leur droite sur la hauteur de la première ; l'indice
$\mathcal{O}_X$ de la première est suivi d'un pâté, et l'indice $X/S$ du
dernier $\underline{\mathrm{Hom}}$ de la troisième est barré d'un trait}

\section*{Problèmes ouverts}

\note{titre de sa main, à l'encre, dans l'angle supérieur droit d'une seconde
chemise brune (page 9), qui ne reçoit pas de \texttt{page}}

\page{10}

\subsection*{Problèmes-types à résoudre}

\note{les lettres a), c) et j) sont encadrées d'un trait épais ; dans la marge
gauche, entourés, des renvois aux paragraphes du plan et à un chapitre VI. Ils
sont donnés ici en marginales, devant l'item qu'ils accompagnent}

\marginal{§ 7}

a) \quad Anneaux de définition pour des structures sans automorphismes

\marginal{§§ 20, 21 ; Chap VI}

b) \quad Hilbert, Picard sans hypothèse projective \quad [Résultats partiels]

\uncertain{Dual} d'un schéma abélien non projectif \quad [OK par Raynaud]

\note{« OK par Raynaud » est encadré ; « Chap VI » est écrit en regard de cette
seconde ligne, le VI repassé au feutre}

\marginal{§§ 3, 4}

c) \quad Passage au quotient sans hypothèse de platitude

\note{double trait horizontal}

\marginal{Mumford}

d) \quad Passage aux quotients par $\mathrm{Gl}_n$, opérant « sans point fixe »

e) \quad Modules de variétés plus ou moins quelconques. --- Cas des schémas
abéliens

\note{d) et e) sont pris dans un crochet épais ; « Mumford », entouré, est relié
à ce crochet}

\marginal{Mumford-Seshadri-Narasimhan \uncertain{pour} fibrés vectoriels ; cas
fini : \ill{}}

f) \quad Classification de fibrés principaux sur $X/Y$ (par exemple) à groupe
affine $G$ sur $Y$, en distinguant les cas
\begin{itemize}
  \item $G$ commutatif
  \item $G$ non commutatif
\end{itemize}

\note{« sur $X/Y$ » est interligné au-dessus de « principaux » avec un trait de
renvoi ; « (par exemple) » est d'une encre plus pâle. La marginale est
entourée et rattachée au crochet qui enserre f). Un « § 17 », entouré, est
relié par un trait à la ligne « $G$ commutatif » et à g)}

\marginal{§ 17}

g) \quad Cas particulier ; $G$ fini :

\marginal{Pb mal posé ?}

h) \quad Cas particulier : \quad $X = \mathbb{P}^1_Y$, \quad $G = \mathrm{Gl}(n)$.

\note{double trait horizontal}

\marginal{Chap VI ; Oui}

i) \quad Les modules formels d'une variété abélienne \struck{\ill{}} sont-ils
simples sur \uncertain{l'anneau de base} ?

\note{« simples » souligné ; « Oui » est entouré, sous « Chap VI » dont le VI est
repassé au feutre}

\note{double trait horizontal}

j) \quad « Th. du cône » dans cas non projectif \ldots

\end{document}
